
\section{Atlante delle matrici del Transformer}
\label{sec:atlante-matrici}

Una difficoltà frequente nella lettura dei risultati consiste nel capire \emph{quali} matrici siano state effettivamente strutturate. RowMask23 non sostituisce un Transformer con un algoritmo diverso: cambia il modo in cui sono rappresentate le grandi proiezioni lineari, mentre l'attenzione causale, i collegamenti residui, le normalizzazioni e la perdita autoregressiva continuano a svolgere i rispettivi ruoli.

La \autoref{fig:atlante-matrici} riassume le dimensioni delle matrici usate dal solo modello addestrato STERA 30M. La matrice dei token è condivisa fra embedding iniziale e testa di uscita. Ogni strato contiene invece quattro proiezioni dell'attenzione e due proiezioni del feed-forward, tutte organizzate in blocchi da $128\times128$.

\begin{figure}[htbp]
\centering
\begin{tikzpicture}[x=1.6mm, y=1.6mm, font=\small,
  cel/.style={draw=steraBlu!70, fill=steraBlu!10, line width=.2pt}]
% E: 393 x 6 blocchi, disegnata compressa (24 righe rappresentative)
\foreach \i in {0,...,5} \foreach \j in {0,...,23} \draw[cel] (\i,-\j) rectangle ++(1,-1);
\node[font=\footnotesize] at (3,-25.6) {$\vdots$};
\foreach \i in {0,...,5} \foreach \j in {27,28} \draw[cel] (\i,-\j) rectangle ++(1,-1);
\node[align=center, anchor=north] at (3,-30) {$E$\\{\footnotesize $50.304\times768$}\\{\footnotesize $393\times6$ blocchi}};
% Q, K, V, O: 6x6
\foreach \m/\x in {Q/12, K/21, V/30, O/39} {
  \foreach \i in {0,...,5} \foreach \j in {0,...,5} \draw[cel] (\x+\i,-\j) rectangle ++(1,-1);
  \node[anchor=north] at (\x+3,-7) {$W_{\m}$};
}
\node[anchor=north, font=\footnotesize, align=center] at (28.5,-10.5) {attenzione: $768\times768$, $6\times6$ blocchi ciascuna};
% FF1 24 x 6, FF2 6 x 24
\foreach \i in {0,...,5} \foreach \j in {0,...,23} \draw[cel] (12+\i,-15-\j) rectangle ++(1,-1);
\node[anchor=north, align=center] at (15,-40) {$W_1$\\{\footnotesize $3.072\times768$}};
\foreach \i in {0,...,23} \foreach \j in {0,...,5} \draw[cel] (22+\i,-15-\j) rectangle ++(1,-1);
\node[anchor=west, align=left] at (46.5,-18) {$W_2$\\{\footnotesize $768\times3.072$}};
\draw[decorate, decoration={brace, amplitude=4pt}, steraGrigio] (11.5,1) -- (46.5,1) node[midway, above=4pt, font=\footnotesize] {ripetute in ciascuno dei 6 strati};
\node[anchor=north west, font=\footnotesize, align=left, text width=40mm] at (22,-25) {ogni quadrato è un blocco\\RowMask23 $128\times128$ con\\3.176 coordinate e una scala};
\end{tikzpicture}

\caption{Inventario visivo delle matrici strutturate di STERA 30M: una matrice condivisa per ingresso e uscita, più quattro proiezioni di attenzione e due del feed-forward in ciascuno dei sei strati. La figura mostra le dimensioni effettivamente impiegate, non uno schema architetturale generico.}
\label{fig:atlante-matrici}
\end{figure}

\subsection{Embedding e testa di uscita: una sola matrice}

Il tokenizzatore dispone di \num{50257} identificativi validi. Per dividere i tensori in blocchi da 128 elementi, la matrice fisica $E$ ha \num{50304} righe: le 47 righe eccedenti non rappresentano token validi. La dimensione dei vettori di embedding è 768. Ne deriva una matrice
\[
E\in\mathbb{R}^{50304\times768},
\qquad
\frac{50304}{128}\cdot\frac{768}{128}=393\cdot6=2358
\]
blocchi RowMask23. Quando il token $u_t$ entra nella rete, si usa la riga $E_{u_t,:}$; dopo i sei strati, lo stesso tensore fornisce i coefficienti per calcolare i punteggi di uscita. Il \emph{weight tying} significa dunque condivisione dei parametri, non la presenza di due tabelle uguali memorizzate separatamente.

Per l'embedding la rappresentazione densa contiene
\[
50304\cdot768=\num{38633472}
\]
valori, mentre le sole coordinate RowMask23 sono
\[
2358\cdot3176=\num{7489008}.
\]
Questi numeri confrontano la componente strutturata con la matrice ricostruita: non comprendono ancora i fattori globali a basso rango né i guadagni per riga e colonna.

\subsection{Le quattro proiezioni dell'attenzione}

Ciascuno dei sei strati contiene le matrici $W_Q$, $W_K$, $W_V$ e $W_O$ di dimensione $768\times768$. Le prime tre trasformano le rappresentazioni in query, chiavi e valori; la quarta riporta il risultato delle teste nello spazio residuo del Transformer. Non si tratta di quattro matrici alternative, ma di quattro operazioni distinte necessarie nello stesso blocco di attenzione.

Una proiezione $768\times768$ contiene $6\cdot6=36$ blocchi. Le quattro proiezioni di uno strato ne richiedono $4\cdot36=144$; in sei strati il totale diventa
\[
6\cdot4\cdot36=\num{864}\quad\text{blocchi}.
\]
L'attenzione del modello usa 16 teste da 48 componenti e alterna tre strati con visibilità sull'intero prefisso e tre con finestra causale di 256 posizioni. Questa scelta riguarda quali token possono interagire tra loro; la parametrizzazione STERA riguarda invece la rappresentazione dei pesi delle proiezioni.

\subsection{Le due proiezioni del feed-forward}

Il feed-forward espande ciascun vettore da 768 a 3072 componenti, applica l'attivazione GELU e lo riporta a 768 componenti. La prima matrice ha dimensione $3072\times768$, dunque $24\cdot6=144$ blocchi; la seconda ha dimensione $768\times3072$, dunque $6\cdot24=144$ blocchi. I due versi sono distinti e devono essere contati entrambi.

Nei sei strati vi sono quindi \num{864} blocchi per l'espansione e \num{864} per la proiezione di ritorno. Sommando attenzione e feed-forward, ogni strato possiede $144+144+144=432$ blocchi. I sei strati, da soli, contengono $6\cdot432=\num{2592}$ blocchi; con la matrice $E$ condivisa il totale è
\[
\num{2358}+\num{2592}=\boxed{\num{4950}\ \text{blocchi RowMask23}}.
\]
La stessa contabilità spiega le 37 matrici del modello: una matrice condivisa più $6\cdot(4+2)=36$ proiezioni interne.

\subsection{Verifica contabile dal blocco al modello}
\label{sec:contabilita-matrici}

La \autoref{tab:atlante-matrici} attraversa la catena completa: numero di matrici, numero di blocchi, coordinate compresse e valori dell'espansione. Per evitare doppi conteggi, $E$ compare una volta sola.

\begin{table}[htbp]
\centering
\small
\begin{tabularx}{\linewidth}{Xrrr}
\toprule
Famiglia di matrici & Matrici & Blocchi & Coordinate RowMask23 \\
\midrule
Embedding e testa condivisi $E$ & 1 & \num{2358} & \num{7489008} \\
Attenzione $Q,K,V,O$ & 24 & \num{864} & \num{2744064} \\
Feed-forward: espansione & 6 & \num{864} & \num{2744064} \\
Feed-forward: ritorno & 6 & \num{864} & \num{2744064} \\
\midrule
\textbf{Totale} & \textbf{37} & \textbf{\num{4950}} & \textbf{\num{15721200}} \\
\bottomrule
\end{tabularx}
\caption{Inventario esatto delle matrici e delle sole coordinate RowMask23 presenti in STERA 30M. Ogni blocco contiene \num{3176} coordinate. Le altre componenti allenabili sono riportate nella \autoref{tab:stera30m-parametri}.}
\label{tab:atlante-matrici}
\end{table}

Gli elementi delle matrici dense corrispondenti sono
\[
\num{38633472}+3\cdot\num{14155776}
=\num{81100800}.
\]
Questo è il numero di elementi delle 37 \emph{matrici principali}, non il numero complessivo di elementi del checkpoint espanso. Aggiungendo posizioni, normalizzazioni e bias, il conteggio totale riportato nell'inventario generale sale a \num{82721280}. Analogamente, le \num{15721200} coordinate dei blocchi non esauriscono i parametri allenabili: per ottenere \num{30346182} occorre includere anche scale, guadagni, fattori $U,V$ di rango 96 e gli altri tensori. La trasparenza di queste due somme è essenziale: chiamare «30M» un modello non equivale a dichiarare che la sua espansione operi con soli 30 milioni di moltiplicatori densi.

\subsection{Come leggere la tabella senza confondere memoria e calcolo}

Le coordinate RowMask23 sono i dati memorizzati; i valori densi sono i coefficienti matematici che la regola di ricostruzione può produrre. I due motori misurati nel capitolo sui tempi di esecuzione rappresentano quindi \emph{lo stesso} STERA 30M, ma organizzano diversamente gli accessi in memoria. Nel motore denso si costruiscono e si mantengono residenti le matrici ricostruite; nel percorso compatto si usano direttamente le coordinate, calcolando i contributi durante il prodotto matrice--vettore.

Non si deve sommare la memoria della matrice densa a quella delle coordinate quando si confrontano i due formati come se fosse una sola rappresentazione obbligatoria. Nella pratica una singola implementazione può avere buffer temporanei e cache che aggiungono memoria al payload persistente: tali costi sono misurati e discussi separatamente. La riduzione teorica dei coefficienti di un singolo blocco non può inoltre essere trasferita indistintamente all'intero Transformer, perché i suoi parametri non sono tutti rappresentati con RowMask23.
